\paragraph{Classical theta sections supply the tangent estimate.}
The arithmetic inputs to the preceding jet argument can be supplied on
the prescribed lattice. The distinction between an even theta level and
an arbitrary positive level is retained in the parameter choice below.

\begin{theorem}\label{geo-classical-tangent-slope}
Let $G/\mathcal O_K$ be a smooth separated group scheme of finite type
with geometrically connected semiabelian fibres and abelian generic
fibre $A$ of dimension $g\ge1$. Let $\mathcal L_0$ be a symmetric
cubical line bundle, rigidified at zero, whose generic fibre $L$
defines a principal polarization. Give it the canonical metric and
give $V=\operatorname{Lie}(G)$ the metric
$\|v\|_\sigma^2=\pi H_\sigma(v,v)$, where $H_\sigma$ is the
Appell--Humbert form of $L$. With ordinary normalized Arakelov degree
$d=\widehat{\deg}_{\rm n}V$, one has
\begin{equation}\label{geo-classical-tangent-twelve}
 \widehat\mu_{\max}(V)
       \le12\log\bigl(2+g+(-d)_+\bigr).
\end{equation}
Here maximal slope is the supremum of the mean degrees of nonzero
integral submodules with inherited metrics. The bound has the same
constant after every finite extension of $K$.
\end{theorem}
\begin{proof}
\emph{The same integral lattice and line.}
The semistable comparison theorem identifies $G$ with the identity
component of the N\'eron model of $A$, and preserves that identity
component under finite extension
\cite[Theorem~4.4 and Corollary~4.5]{conrad-semistable}.
Thus the identity component of a Moret--Bailly model over an extension
$K'$ is $G_{\mathcal O_{K'}}$; its tangent module is precisely
$V\otimes\mathcal O_{K'}$.

The Moret--Bailly extension of the specified $L$ restricts to
$\mathcal L_0$ on this component. Indeed the divisor of the generic
rigidified isomorphism is vertical. Locally over a discrete valuation
ring, the total space is regular and its special fibre is integral,
so that divisor is a multiple of the special fibre. Pullback at zero
shows that its multiplicity is zero. Normality extends the isomorphism,
and generic density makes it unique. Its canonical metrics agree:
their logarithmic ratio is harmonic on the compact complex torus and
vanishes at the zero rigidification. No integral index or metric
change is introduced.

\emph{One integral theta section.}
In symplectic coordinates with period matrix $\tau$, the squared norm
on the dual tangent determinant is
\[
 \det(\operatorname{Im}\tau)/\pi^g
  =\frac{i^{g^2}}{(2\pi)^g}\int_A
        dz_1\wedge\cdots\wedge dz_g\wedge
        d\bar z_1\wedge\cdots\wedge d\bar z_g.
\]
At finite places this is the invariant Hodge line of the same model.
Consequently the Faltings height in this convention is $h_F(A)=-d$.
Write $h=(-d)_+$ and $Y=2+g+h$.

For a Moret--Bailly model $\pi:\mathcal A\to\operatorname{Spec}
\mathcal O_{K'}$, the rank-one module $\pi_*\mathcal L$, with the
$L^2$ norm for Haar probability measure, has degree
$-h_F(A)/2-g\log(2\pi)/4$. It is preserved isometrically under
change of model and finite extension
\cite[Definition~2.1 and Theorem~3.4(i),(iii)--(v)]{pazuki-theta-faltings}.
Principalize this rank-one module by the ideal argument used above
for the absolute Siegel lemma. An integral generator $\theta$ then
satisfies
\begin{equation}\label{geo-theta-ltwo-supplier}
 \operatorname{av}\log\|\theta\|_2
          \le h/2+g\log(2\pi)/4.
\end{equation}
Restriction gives an integral section of the prescribed line bundle.

For every positive even integer $m$, extend the field to contain
$A[m^2]$ and use a model containing these torsion sections. The normalized
theta-group maps $[m]^*L\simeq L^{m^2}$ and
$t_a^*L^{m^2}\simeq L^{m^2}$ extend integrally and isometrically
\cite[Propositions~2.5 and~5.1]{pazuki-theta-faltings}.
Representatives $a\in A[m^2]/A[m]$ produce $m^{2g}$ integral
sections $s_a$ spanning $H^0(A,L^{m^2})$, with
\[
 \|s_a(z)\|=\|\theta([m](z+a))\|,
 \qquad \|s_a\|_{\infty,\sigma}=\|\theta\|_{\infty,\sigma}.
\]
The first section remains integral on each new model by model
independence. The integral spanning assertion is only generic spanning;
an integral basis of the whole section module is unnecessary.

\emph{A uniform supremum bound.}
Let $\rho_\sigma$ be the shortest nonzero period in the metric
$H_\sigma$, and put $r_\sigma=\min(1,\rho_\sigma/2)$.
The principal period lattice has real $H_\sigma$-covolume one.
In a Gaussian holomorphic frame centred at any point, the squared
norm is $|f(w)|^2e^{-\pi H_\sigma(w,w)}$. The ball of radius
$r_\sigma$ embeds in the torus. Its volume in Haar probability units
is $\pi^g r_\sigma^{2g}/g!$, so the holomorphic mean-value inequality gives
\[
 \frac{\|\theta\|_{\infty,\sigma}^2}{\|\theta\|_{2,\sigma}^2}
       \le\frac{g!e^{\pi r_\sigma^2}}{\pi^g r_\sigma^{2g}}.
\]
Autissier's matrix lemma uses the height
$h_{\rm Fa}=h_F-g\log\pi/2$. Its corollary at $\varepsilon=1/2$
gives exactly
\cite[Corollary~1.4 and the height convention in Section~2]{autissier-matrix}
\[
 \operatorname{av}\rho_\sigma^{-2}
 \le\frac{12}{\pi}h_F+\frac{6g}{\pi}\log(4\pi)
 \le4h+6g,\qquad
 \operatorname{av}r_\sigma^{-2}\le24Y.
\]
Take logarithms, average, and use Jensen's inequality together with
$\log(g!)\le g\log g$. Combining with
\eqref{geo-theta-ltwo-supplier} yields, for
$M_\sigma=\|\theta\|_{\infty,\sigma}$,
\begin{equation}\label{geo-theta-sup-supplier}
 U:=\operatorname{av}\log M_\sigma
                 \le h/2+5g+g\log Y.
\end{equation}
For the constant, it suffices that
$2+\tfrac12\log24+\tfrac14\log(2\pi)<5$.
All these averages are unchanged under finite extension.

\emph{Even level and the slope.}
Set
\[
 Z=1+11g+h+2g\log Y,\quad a=\lceil\sqrt Z\rceil,
 \quad m=6a,\quad k=12a^2,\quad n=m^2=3k.
\]
Then $12Z\le k\le48Z$ and $Z\le13Y\log Y$. The level $m$
is even, so the preceding integral sections apply. In
\eqref{geo-exact-jet-cost} their number is $n^g$, giving every
nonzero integral tangent vector the bound
\[
 -\operatorname{av}\log\|x\|
 \le\log k+\tfrac12\log(3e)+\frac{g\log(3k)+2U}{k}
 \le\frac{13}{12}\log(48Z)+\frac7{12}\log3+\frac12.
\]
Here $g+2U_+\le Z$, $Z/k\le1/12$, and $\log(3k)>1$.
Thus the line bound has coefficient
\[
 C_{\rm line}=\frac{13}{6}+\frac{13}{12}\log624
                              +\frac7{12}\log3+\frac12
\]
in front of $\log Y$. Apply the already established absolute Siegel
comparison, after an isometric finite extension, to each nonzero
integral submodule. Its additional coefficient is at most $1/2$.
Since $\log624<7$ and $\log3<2$,
$C_{\rm line}+1/2<143/12<12$, proving
\eqref{geo-classical-tangent-twelve}.
\end{proof}

For selected submodules of total rank $R$, the theorem supplies
$T\le12R\log(2+g+(-d)_+)$ on these same lattices. In particular
Proposition~\ref{geo-selected-tangent-sum} applies with $C_s=12$.
If the abelian rank is zero, the selected degree sum is zero without
using theta sections. The flat determinant, point-height, cusp and
finite section-degree comparisons are separate geometric requirements.
